\section{Discussion} \label{sec:discussion}

\subsection{Key Findings}

The experimental results reveal several important observations about CPS security evaluation:

\textbf{Pattern-based defences are fragile under adversarial scheduling.}
The Markov defence achieved 52.4\% block rate against static attacks (E2) but only 44.2\% (mean across 20 seeds) against the same attacks delivered with randomised timing and intensity (E3), with F1 dropping from 0.982 to 0.567.
This degradation demonstrates that attackers who vary their behaviour, even without changing the attack type, can reduce defence effectiveness.
This finding aligns with Sommer and Paxson's observation that closed-world evaluation overstates IDS performance~\cite{sommer2010closedworld}.

\textbf{Multi-feature defences recover under adversarial scheduling.}
The Isolation Forest defence (E4), using all 11 features, achieved an 81.2\% mean block rate and F1 = 0.822 against the same randomised attacks that degraded the Markov defence to F1 = 0.567.
The improvement is statistically significant ($p < 0.001$ for recall, F1, and block rate).
This confirms that combining rate, entropy, source diversity, and sequential features provides greater resilience than relying on a single signal (transition probability).

\textbf{Detection is necessary but insufficient.}
E1 achieved 99.2\% recall (the detector correctly identified virtually all attack messages) yet the mission impact ratio reached 62.5\%.
High detection accuracy without active blocking provides situational awareness but does not protect the mission.
This gap between detection and protection shows that evaluating \emph{active} defence mechanisms alongside passive detection is necessary.

\textbf{Mission-level metrics reveal what classification metrics hide.}
Even when the Markov defence's mean F1 drops from 0.982 (E2) to 0.567 (E3, 20-seed mean), the mission impact ratio shifts from 0.139 to 0.978, a far larger proportional change.
This confirms that traditional classification metrics alone are insufficient for CPS security evaluation. Mission-aware metrics are essential for understanding real-world impact.

\textbf{Multi-feature defences outperform simple rate thresholds.}
The rate-threshold baseline, which monitors only total message rate, can detect volumetric floods but fails to identify low-rate spoofing or injection attacks that stay below the rate threshold.
Both the Markov and Isolation Forest defences achieve higher precision than the rate baseline, confirming that learned models add measurable value in distinguishing attack traffic from normal traffic.
The rate baseline achieves higher recall and block rate through aggressive blanket blocking, but at the cost of also blocking legitimate messages.
This comparison addresses the concern that a trivial approach might suffice.

\subsection{The Role of the Orchestrator}

A central contribution of this work is the orchestrator component.
Without automated orchestration, reproducing the E3 experiment would require manually scripting:
\begin{itemize}
    \item Shuffled attack ordering with specific random seed.
    \item Per-segment rate sampling across 5--13 sub-episodes per attack.
    \item Micro-pause insertion between segments.
    \item Synchronised timing of detector, defence, and logging.
    \item Consistent metric computation across configurations.
\end{itemize}

The orchestrator reduces this complexity to a single command-line invocation, making systematic parameter sweeps and fair cross-defence comparison feasible.
The \texttt{batch\_runner.py} module further automates multi-configuration evaluation through YAML-specified sweeps.

\subsection{Limitations}

\textbf{Simulated traffic only.}
The current implementation uses a MAVLink SITL substitute rather than a physical UAV or a high-fidelity simulator (AirSim, Gazebo).
While the MAVLink message types and protocol semantics are authentic, the traffic patterns may not capture all real-world phenomena (sensor noise, GPS multipath, network jitter).

\textbf{Single communication channel.}
All traffic flows through a single UDP port.
Real autonomous systems may use multiple communication channels (telemetry, video, command uplink) with different vulnerability profiles.

\textbf{Baseline training duration.}
The 5-second baseline phase produces approximately 69 training samples, which limits the Isolation Forest's ability to model rare-but-normal message sequences.
Longer baselines may improve defence precision at the cost of extended setup time.

\textbf{Threshold sensitivity.}
All three defences depend on manually chosen thresholds ($\tau_M = 0.05$ for Markov, $\tau_A = 0.5$ for Isolation Forest, $k = 2$ for rate-threshold).
Systematic threshold optimisation (e.g., via ROC analysis on a held-out validation set) is left for future work.

\subsection{Future Work}

\begin{enumerate}
    \item \textbf{High-fidelity simulation integration:} Connect the framework to AirSim~\cite{shah2018airsim} or Gazebo~\cite{koenig2004gazebo}-based UAV missions to measure physical-level mission impact (trajectory deviation, collision avoidance degradation).

    \item \textbf{Adaptive threshold learning:} Implement online threshold adjustment based on observed false-positive rates during the first minutes of monitoring.

    \item \textbf{Additional attack types:} Extend the catalogue with GPS drift injection, camera frame manipulation, and timing-based attacks (delay, reorder).

    \item \textbf{Multi-platform evaluation:} Evaluate the framework on different autonomous platforms (ground vehicles, marine USVs) to validate generalisation.

    \item \textbf{Containerisation:} Package the framework as a Docker container for one-command reproducibility, fulfilling the original project proposal's containerisation objective.
\end{enumerate}
